\chapter{Растворимость: сколько водорода хочет выпасть}\label{ch:solubility}

\section{Равновесие раствора и гидрида}

Растворённый водород сидит в междоузлиях решётки циркония. Если раствор разбавленный, его
химический потенциал (свободная энергия на моль водорода) равен
\begin{equation}
\mu_{\text{р-р}}=\mu_0(T)+RT\ln c ,
\end{equation}
где $c$ --- концентрация. У водорода в гидриде потенциал $\mu_{\text{гидр}}(T)$ от $c$ не зависит:
гидрид --- отдельная фаза постоянного состава. Водород переходит туда, где потенциал ниже.
Равновесие --- когда потенциалы равны:
\begin{equation}
c_{\text{eq}}(T)=\exp\frac{\mu_{\text{гидр}}-\mu_0}{RT}=A\,e^{-Q/RT}.
\end{equation}
Это и есть \emph{линия растворимости}: закон Аррениуса с «теплотой растворения» $Q$. В
металловедении её записывают через десятичный логарифм:

\begin{keyf}[Линии растворимости]
\[
\lg C\,[\text{ppm}]=b-k\,\frac{1000}{T\,[\text{K}]},\qquad Q=\ln10\cdot1000\,k\,R .
\]
Э635 (Плясов и др. 2023, ДСК): TSSD --- $b=5.56$, $k=2.24$; TSSP --- $b=5.08$, $k=1.68$.
\end{keyf}

\section{Две линии и гистерезис}

\fig{F03_Solubility}{Цикл опыта на диаграмме «температура--водород в растворе» (Э635, 178 ppm).
При нагреве гидриды растворяются по TSSD, при охлаждении раствор пересыщается и гидрид начинает
выпадать только на TSSP, дальше --- вдоль неё. Между линиями --- гистерезис.}

У циркония растворение и выпадение идут по \emph{разным} линиям (рис.~\ref{fig:F03_Solubility}).
\begin{itemize}
\item \textbf{TSSD} (terminal solid solubility for dissolution) --- при нагреве гидрид растворяется,
  пока концентрация в растворе не дойдёт до этой линии. Это почти равновесие, и она одна и та же
  у Э635, Э110, Zircaloy, M5 и ZIRLO в пределах погрешности.
\item \textbf{TSSP} (for precipitation) --- при охлаждении гидрид начинает выпадать, только когда
  раствор пересыщен до этой линии. Она \emph{кинетическая}: зависит от скорости охлаждения и
  от $T_{\max}$.
\end{itemize}
Почему нужен запас? Гидрид на 15\% объёмнее металла, и зародышу приходится раздвигать решётку
(гл.~\ref{ch:misfit}). На это уходит энергия, которую должна оплатить «лишняя» химическая
движущая сила. В модели TSSD --- равновесная линия (до неё растут уже выросшие пластинки), а
TSSP --- линия, на которой зарождение становится заметным.

\paragraph{Цикл опыта.} Нагрели до $T_{\max}$: в раствор ушло $\min(H,\ \TSSD(T_{\max}))$.
Если $T_{\max}$ ниже температуры полного растворения, часть гидридов осталась --- это
\emph{неполное растворение}, и оставшиеся (окружные) пластинки потом служат затравкой.
Охлаждаем: концентрация стоит на месте, пока не встретит TSSP, затем раствор следует за TSSP вниз.

\begin{exercise}
Для Э635 и $H=178$~ppm найдите: (а) TSSD при 400\,$^\circ$C; (б) температуру полного
растворения; (в) температуру, при которой начнётся выпадение при охлаждении; (г) теплоту $Q$
для TSSD и для TSSP.
\end{exercise}

\section{Движущая сила в мегапаскалях}

Модель сравнивает химию с механикой, поэтому движущую силу удобно мерить в единицах энергии
на объём гидрида, то есть в мегапаскалях. Моль водорода, перешедший из раствора с концентрацией
$c$ в гидрид при равновесной концентрации $c_{\text{eq}}$, выигрывает $RT\ln(c/c_{\text{eq}})$.
В кубометре гидрида $n_H\approx1.01\cdot10^5$ моль водорода (плотность ZrH$_{1.66}$ 5.65~г/см$^3$,
молярная масса 92.9~г/моль, 1.66 атома H на атом Zr). Значит, выигрыш на единицу объёма
гидрида
\begin{equation}
\Delta_{\text{хим}}=n_H\,RT\ln\frac{c}{c_{\text{eq}}}.
\end{equation}
Множитель $n_H R=0.84$~МПа/К, при 600~K это $\approx500$~МПа на единицу логарифма: пересыщение на
10\% уже дают $\approx50$~МПа. Химия --- сильный игрок.

В модели движущая сила отсчитывается от TSSP:
\begin{equation}
\Delta=\Delta_0+n_H RT\ln\frac{c}{c_{\TSSP}(T)}+\dots
\label{eq:delta_chem}
\end{equation}
На линии TSSP она равна $\Delta_0$ --- это «сколько нужно для заметного зарождения» (15~МПа,
гл.~\ref{ch:nucleation}). Многоточие --- механические добавки, о них дальше.

\section{Градусы и мегапаскали --- один язык}

Многие опыты дают сдвиг \emph{температуры} выпадения: например, нагрузка сдвигает начало выпадения
на $0.08\,^\circ$C/МПа (Vizcaíno 2014), а зёрна разной ориентации выпадают с разницей в
несколько градусов. Как перевести это в движущую силу? Сдвинуть начало выпадения на $\delta T$
вверх --- то же, что заменить $c_{\TSSP}(T)$ на $c_{\TSSP}(T+\delta T)$. Так как
$\ln c_{\TSSP}=\text{const}-Q/RT$,
\[
n_HRT\left[\ln c_{\TSSP}(T+\delta T)-\ln c_{\TSSP}(T)\right]\approx n_HRT\cdot\frac{Q}{RT^2}\,\delta T
=n_H\frac{Q}{T}\,\delta T .
\]

\begin{keyf}[Перевод градусов в мегапаскали]
\[
\frac{\partial\Delta}{\partial T}=n_H\,\frac{Q}{T}\approx\frac{1.01\cdot10^5\cdot32.2\cdot10^3}{603}\ \text{Па/К}
\approx5.4\ \text{МПа}/^\circ\text{C}\quad(\text{TSSP Э635, }330\,^\circ\text{C}).
\]
\end{keyf}

\fig[0.85]{F04_Scales}{Один и тот же сдвиг --- в градусах (сверху) и в мегапаскалях движущей
силы (снизу). Точки --- величины, которые встретятся в книге: фора границы зерна 1.5\,$^\circ$C
(гл.~\ref{ch:nucleation}), фора зёрен 5\,$^\circ$C у Zr--Nb и 10.9--12.4\,$^\circ$C у Zircaloy-4.}

\begin{idea}
Градус --- это много. $1\,^\circ$C сдвига линии выпадения стоит $\approx5$~МПа движущей силы,
а упругая работа окружного напряжения 100~МПа на несоответствии гидрида даёт всего
$\approx3$~МПа (гл.~\ref{ch:misfit}). Поэтому то, \emph{какое семейство зёрен начинает выпадать
первым}, оказывается важнее тонкостей поля напряжений --- это и есть механизм порога
(гл.~\ref{ch:nucleation}).
\end{idea}

\begin{exercise}
При 300\,$^\circ$C найдите для Э635 отношение $c_{\TSSP}/c_{\TSSD}$ и «цену гистерезиса»
$n_HRT\ln(c_{\TSSP}/c_{\TSSD})$ в МПа. Сравните её с упругой энергией несоответствия на единицу
объёма $\sim E\eps^{*2}$ при $E=90$~ГПа и $\eps^*\approx0.06$. Что это говорит о природе гистерезиса?
\end{exercise}

\begin{exercise}
Повторите вывод $\partial\Delta/\partial T=n_HQ/T$ и посчитайте его при 250 и 400\,$^\circ$C.
Насколько сильно «курс обмена» градусов на мегапаскали меняется в интервале выпадения?
\end{exercise}

\begin{codebox}
\textbf{Где в коде.} \code{thermo.py}: линии \code{TSS}, \code{c\_line}, \code{T\_line}, ход
охлаждения \code{Cooling}; \code{ca\_kinetic.py}: \code{N\_H\_MPa\_per\_K} $=n_HR$ и член
$n_HRT\ln(c/c_{\TSSP})$ в движущей силе.
\end{codebox}
